\topsection[sec:design]{Solution design}

This section is the summary of phase~3. Three documents hold the detail: the architecture with its six
diagrams (\repofile{docs/05_architecture.md}), the functional description
(\repofile{docs/06_functional_description.md}) and the technical description
(\repofile{docs/07_technical_description.md}).

\subsection{Overview}
\sys\ attests four stages of an AI pipeline (Figure~\ref{fig:arch}). At each stage a role signs a manifest
that describes the artifact and names its parents by the digest of their own manifests. A verifier checks an
artifact together with its whole ancestry, a signed trust log says which keys count and what has been revoked,
and a deployment gateway refuses to load a model whose chain does not verify.

\begin{figure}[!htb]
\centering
\begin{tikzpicture}[x=1cm,y=1cm,font=\sffamily\fontsize{8}{10}\selectfont,
  stage/.style={draw=cgnavy,line width=0.7pt,rounded corners=3pt,minimum width=3.05cm,minimum height=1.65cm,align=center,fill=white,inner sep=2pt},
  side/.style={draw=cgmuted,line width=0.6pt,rounded corners=3pt,minimum height=1.35cm,align=center,fill=cgband,inner sep=3pt},
  lnk/.style={-{Latex[length=2.1mm,width=1.7mm]},line width=0.8pt,color=cgteal},
  rep/.style={-{Latex[length=2mm,width=1.6mm]},line width=0.7pt,color=cgmuted,densely dashed},
  role/.style={font=\ttfamily\fontsize{6.7}{8}\selectfont,text=cgmuted}]
\node[stage] (s1) at (0,0) {\textbf{\color{cgnavy}Datasets}\\stage 1\\[1pt]{\ttfamily\fontsize{6.7}{8}\selectfont\color{cgmuted}data\_steward}};
\node[stage] (s2) at (3.98,0) {\textbf{\color{cgnavy}Base model}\\stage 2\\[1pt]{\ttfamily\fontsize{6.7}{8}\selectfont\color{cgmuted}training\_pipeline}};
\node[stage] (s3) at (7.96,0) {\textbf{\color{cgnavy}Fine-tuned model}\\stage 3\\[1pt]{\ttfamily\fontsize{6.7}{8}\selectfont\color{cgmuted}training\_pipeline}};
\node[stage] (s4) at (11.94,0) {\textbf{\color{cgnavy}Deployment record}\\stage 4, the gateway\\[1pt]{\ttfamily\fontsize{6.7}{8}\selectfont\color{cgmuted}deployment\_gateway}};
\node[side,minimum width=3.05cm] (ext) at (3.98,-2.35) {\textbf{\color{cgnavy}External base model}\\{\ttfamily\fontsize{6.7}{8}\selectfont\color{cgmuted}model\_publisher}\\we hold its public key only};
\node[side,minimum width=4.85cm] (log) at (11.04,-2.35) {\textbf{\color{cgnavy}Trust log}\\signed, hash-chained records; pinned head\\{\ttfamily\fontsize{6.7}{8}\selectfont\color{cgmuted}release\_manager} (root key)};
\draw[lnk] (s1) -- (s2);
\draw[lnk] (s2) -- (s3);
\draw[lnk] (s3) -- (s4);
\draw[lnk] (s1.south) -- ++(0,-0.42) -| ([xshift=-9mm]s3.south);
\draw[lnk] (ext.east) -| (s3.south);
\draw[rep] (log.north -| s4.south) -- (s4.south);
\node[font=\sffamily\fontsize{7.1}{8.6}\selectfont,text=cgmuted,anchor=north] at (1.99,-1.3) {fine-tuning dataset};
\node[font=\sffamily\fontsize{7.1}{8.6}\selectfont,text=cgmuted,anchor=east,align=right] at (11.8,-1.25) {replayed at\\every verification};
\end{tikzpicture}
\caption{The four attested stages and the role that signs each one. A solid arrow means that the signed
payload of the child embeds the manifest digest of the parent. An internal base model must have a dataset
parent. An external base model has no internal lineage and may be signed only by a registered publisher key.}
\label{fig:arch}
\end{figure}

\subsection{Signed manifests and the hash-linked chain}
Every artifact receives a manifest body $m$: its type and pipeline stage, the SHA-256 and the size of the file,
for CSV datasets the Merkle root over the hashes of its rows, the list of its parents, the provenance (training
configuration and environment), the results of the content checks made before signing, and the identity, role
and key identifier of the signer. The body is serialised in a canonical JSON form (sorted keys, no
insignificant whitespace) and signed as a whole with Ed25519:
\begin{equation}
\sigma = \mathrm{Sign}_{k}\bigl(\mathrm{canon}(m)\bigr),\qquad
d(m) = \mathrm{SHA256}\bigl(\mathrm{canon}(m)\bigr).
\label{eq:sign}
\end{equation}
Signing the whole body, and not the file hash alone, prevents an attacker from keeping a valid pair of hash
and signature while changing the stage, the parents or the role around it. The manifest digest $d(m)$ is the
identity of the artifact in the provenance graph. A child lists each parent as the triple of its identifier,
its manifest digest and the relation, inside its own signed body (Figure~\ref{fig:chain}). The digest of a
child therefore commits to the digest of every ancestor, as a Git commit commits to its parents. Rewriting an
ancestor, even by signing it again with the legitimate key, changes its digest and breaks every descendant.

\begin{figure}[!htb]
\centering
\begin{tikzpicture}[x=1cm,y=1cm,font=\sffamily\fontsize{7.6}{9.6}\selectfont,
  head/.style={fill=cgnavy,text=white,minimum width=4.5cm,minimum height=0.56cm,font=\sffamily\bfseries\fontsize{8}{10}\selectfont,inner sep=0pt},
  body/.style={draw=cgnavy,line width=0.6pt,minimum width=4.5cm,minimum height=2.3cm,inner sep=0pt,fill=white},
  key/.style={anchor=west,font=\ttfamily\fontsize{6.5}{8.2}\selectfont,text=cgmuted},
  val/.style={anchor=west,font=\ttfamily\fontsize{6.5}{8.2}\selectfont,text=cgink},
  dig/.style={draw=cgteal,line width=0.7pt,rounded corners=2pt,fill=white,font=\ttfamily\fontsize{7}{8.6}\selectfont,text=cgteal,inner xsep=4pt,inner ysep=2.2pt},
  lnk/.style={-{Latex[length=2mm,width=1.6mm]},line width=0.8pt,color=cgteal}]
\foreach \x/\name/\ttl/\pa/\who/\ek/\ev in {
  0/a/{Manifest of the dataset}/{none}/{data\_steward}/{merkle.root}/{root of the rows},
  5.2/b/{Manifest of the base model}/{d(dataset)}/{training\_pipeline}/{provenance}/{epochs, seed},
  10.4/c/{Manifest of the fine-tuned model}/{d(base), d(data)}/{training\_pipeline}/{provenance}/{epochs, seed}}{
  \node[body] (\name) at (\x,0) {};
  \node[head] at (\x,1.43) {\ttl};
  \node[key] at (\x-2.2,0.82) {subject.sha256};  \node[val] at (\x-0.2,0.82) {file hash};
  \node[key] at (\x-2.2,0.41) {\ek};             \node[val] at (\x-0.2,0.41) {\ev};
  \node[key] at (\x-2.2,0.0) {parents};          \node[val] at (\x-0.2,0.0) {\pa};
  \node[key] at (\x-2.2,-0.41) {signer.role};    \node[val] at (\x-0.2,-0.41) {\who};
  \node[key] at (\x-2.2,-0.82) {signature};      \node[val] at (\x-0.2,-0.82) {whole body};
  \node[dig] (d\name) at (\x,-1.72) {d(m) = SHA256(canon(m))};
}
\draw[lnk] (2.95,0) -- (2.6,0) |- (da.east);
\draw[lnk] (8.15,0) -- (7.8,0) |- (db.east);
\end{tikzpicture}
\caption{Each manifest is signed as a whole. Its \code{parents} field holds the manifest digest of each
parent (arrows). The fine-tuned model has two parents: the base model and the fine-tuning dataset, whose
manifest is not drawn.}
\label{fig:chain}
\end{figure}

For datasets, a Merkle tree over the hashes of the lines (the construction of RFC~6962 \cite{rfc6962}) lets the
verifier report \emph{which} rows changed, authenticated by the signed root. The hashes of the leaves are kept
in a side file that is trusted only if it reproduces that root.

\subsection{Roles and policy}
A valid signature only proves who signed. A policy file states who is \emph{allowed} to sign what and how
artifacts may be connected (Table~\ref{tab:policy}). The verifier enforces it, so a stolen data-steward key
cannot sign a model, and a model without a dataset parent is rejected even when it is correctly signed. The
training key may not sign the external-import stage, so an internal model cannot be relabelled as external to
escape the dataset rule.

\begin{table}[H]
\caption{The supply-chain policy (\code{config/policy.yaml}): one signing role per stage, and the parents that each artifact type requires.}
\label{tab:policy}
\tablefont
\begin{tabularx}{\textwidth}{P{3.3cm}P{3.3cm}P{3.2cm}L}
\toprule
\headrow\hd{Stage} & \hd{Artifact type} & \hd{Signing role} & \hd{Required parents} \\
\midrule
1\enskip Dataset & dataset & \code{data\_steward} & none \\
2\enskip Training & base model (internal) & \code{training\_pipeline} & a dataset \\
2\enskip External import & external base model & \code{model\_publisher} & none: provenance ends at the publisher \\
3\enskip Fine-tuning & fine-tuned model & \code{training\_pipeline} & a dataset and one base model of either kind \\
4\enskip Deployment & deployment record & \code{deployment\_gateway} & a fine-tuned model \\
\bottomrule
\end{tabularx}
\end{table}

\subsection{Verification: ten checks and the status}
Verifying an artifact runs ten checks and recurses into its parents (Table~\ref{tab:checks}). All ten always
run, so a report is complete. The status is derived and never set by hand (Figure~\ref{fig:ladder}):
\st{FAILED} if any check of integrity, trust, provenance, revocation or policy fails; \st{WARNING} if integrity
holds but a content check flagged the data here or upstream; \st{VERIFIED} otherwise. An external model that
arrives unsigned or malformed is \st{REFUSED} at import and nothing is stored. \st{VERIFIED} means \emph{exactly
what trusted, unrevoked signers in authorised roles attested, through an intact ancestry}. The system never
upgrades it to ``safe''.

{\tablefont
\begin{xltabular}{\textwidth}{P{0.75cm}P{3.2cm}L}
\caption{The ten checks of a verification.}\label{tab:checks}\\
\toprule
\headrow\hd{} & \hd{Check} & \hd{Fails when} \\
\midrule
\endfirsthead
\toprule
\headrow\hd{} & \hd{Check} & \hd{Fails when} \\
\midrule
\endhead
\bottomrule
\endfoot
V0 & Envelope & the manifest is missing, malformed, or stored under another identifier \\
V1 & File hash & the SHA-256 or the size differs from the signed values \\
V2 & Dataset Merkle root & rows changed (the report lists which lines) \\
V3 & Ed25519 signature & it is invalid over the canonical body, or the key is not in the signed trust log \\
V4 & Signer trust & the key is revoked, the role or the identity does not match, the role is not allowed for the stage, or the index and the log disagree \\
V5 & Parent links & a parent manifest is missing, or its recomputed digest differs from the embedded one \\
V6 & Ancestry & any ancestor fails (recursive) \\
V7 & Revocation & the artifact or an ancestor is revoked in the signed log, its index row was deleted, or the index flag was cleared against the log \\
V8 & Policy & wrong type for the stage, missing or forbidden parent types, or a parent from a later stage \\
V9 & Content flags & a signed content check flagged the data here or upstream: \st{WARNING} only, never \st{FAILED} \\
\end{xltabular}}

\begin{figure}[!htb]
\centering
\begin{tikzpicture}[x=1cm,y=1cm,font=\sffamily\fontsize{7.5}{9.3}\selectfont,
  gate/.style={draw=cgnavy,line width=0.65pt,rounded corners=2.5pt,minimum width=1.95cm,minimum height=1.3cm,align=center,fill=white,inner sep=1.5pt},
  go/.style={-{Latex[length=1.9mm,width=1.5mm]},line width=0.75pt,color=cgnavy},
  no/.style={-{Latex[length=1.8mm,width=1.4mm]},line width=0.7pt,color=cgmuted}]
\foreach \i/\x/\ttl/\what in {
  1/0/{V0 to V2}/{manifest, file\\hash, rows},
  2/2.25/{V3, V4}/{signature,\\signer and role},
  3/4.5/{V5, V6}/{parent links,\\ancestry},
  4/6.75/{V7}/{revocation},
  5/9.0/{V8}/{policy},
  6/11.25/{V9}/{content flags}}{
  \node[gate] (g\i) at (\x,0) {\textbf{\color{cgnavy}\ttl}\\\what};
}
\foreach \a/\b in {1/2,2/3,3/4,4/5,5/6} \draw[go] (g\a) -- (g\b);
\node[fill=cgteal,text=white,rounded corners=2pt,font=\sffamily\bfseries\fontsize{7.2}{8}\selectfont,inner xsep=5pt,inner ysep=3.5pt] (ok) at (13.3,0) {VERIFIED};
\draw[go] (g6) -- (ok);
\node[fill=cgfailed,text=white,rounded corners=2pt,font=\sffamily\bfseries\fontsize{7.2}{8}\selectfont,minimum width=10.95cm,inner ysep=3.5pt] (fail) at (4.5,-1.55) {FAILED: the gateway blocks};
\foreach \i in {1,2,3,4,5} \draw[no] (g\i.south) -- (g\i.south |- fail.north);
\node[fill=cgwarning,text=white,rounded corners=2pt,font=\sffamily\bfseries\fontsize{7.2}{8}\selectfont,inner xsep=5pt,inner ysep=3.5pt,align=center] (warn) at (12.05,-1.55) {WARNING: blocked by default};
\draw[no] (g6.south) -- (g6.south |- warn.north);
\end{tikzpicture}
\caption{How the status is derived. A failure in V0 to V8 gives \st{FAILED}. A content flag alone gives
\st{WARNING}, which the gateway blocks unless an operator overrides it, and the override is logged.}
\label{fig:ladder}
\end{figure}

\subsection{The signed trust log}
\label{sec:trustlog}
In version 0.3 of the system (October 2026) the trust store, which says which public key belongs to which
identity and role, and the revocation flags lived in the SQLite index. Our own revocation experiment showed the
consequence: an attacker able to write that file could clear a \code{revoked} flag, and the withdrawn model
verified again. Since version 0.4, trust and revocation are records of an append-only log
(Figure~\ref{fig:log}). Record $i$ carries a sequence number, the hash $h_{i-1}$ of the previous record, a
timestamp, an action (\code{trust\_key}, \code{revoke\_key}, \code{revoke\_artifact}, \code{rotate\_root}) and
its subject, and is signed by the root key of the \code{release\_manager} role. A pin file in the trusted keys
directory holds the root public key and the pair $(n, h_n)$ of the last record.

\begin{figure}[!htb]
\centering
\begin{tikzpicture}[x=1cm,y=1cm,font=\sffamily\fontsize{7.3}{9}\selectfont,
  rec/.style={draw=cgnavy,line width=0.6pt,rounded corners=2.5pt,minimum width=2.52cm,minimum height=1.5cm,align=center,fill=white,inner sep=1.5pt},
  back/.style={-{Latex[length=1.8mm,width=1.4mm]},line width=0.7pt,color=cgteal},
  pinl/.style={-{Latex[length=1.9mm,width=1.5mm]},line width=0.7pt,color=cgnavy,densely dashed}]
\foreach \i/\x/\act/\subj in {
  1/0/{trust\_key}/{release\_manager},
  2/2.9/{trust\_key}/{data\_steward},
  3/5.8/{trust\_key}/{training\_pipeline},
  4/8.7/{trust\_key}/{deployment\_gateway},
  5/11.6/{revoke\_artifact}/{base model}}{
  \node[rec] (r\i) at (\x,0) {\textbf{\color{cgnavy}record \i}\\{\ttfamily\fontsize{6.5}{8}\selectfont \act}\\{\fontsize{6.9}{8.4}\selectfont\color{cgmuted}\subj}\\{\fontsize{6.6}{8}\selectfont\color{cgmuted}signed by the root key}};
}
\foreach \a/\b in {2/1,3/2,4/3,5/4} \draw[back] (r\a) -- (r\b);
\node[draw=cgnavy,line width=0.9pt,rounded corners=2.5pt,fill=cgband,align=center,inner xsep=6pt,inner ysep=4pt] (pin) at (10.25,-1.8) {\textbf{\color{cgnavy}Pin file} (trusted keys directory)\\root public key, last record: $(5, h_5)$};
\draw[pinl] (pin) -- (r5);
\node[anchor=west,align=left,text=cgmuted,text width=8.3cm] at (-1.26,-1.8) {Each record embeds the hash of the previous one. The verifier replays the log at every verification: every signature must verify under the root key, every record must chain, and the log must end exactly at the pinned head.};
\end{tikzpicture}
\caption{The trust log after a revocation: four records that trust the role keys and one that revokes an
artifact. There is no action that undoes a revocation; a withdrawn artifact is replaced by a new version.}
\label{fig:log}
\end{figure}

A truncated log (a rollback), an edited record, a deleted or substituted log, and a key that is present in the
index but absent from the log all make V3, V4 and V7 fail closed. An index that disagrees with the log is
reported as tampering. The structure is a single-writer hash-chained log, small enough to need no Merkle tree,
with the root pinning and rollback protection of The Update Framework \cite{tuf2010}.

\subsection{Deployment gateway and external publishers}
The gateway is the last line of defence. It verifies the full chain at deployment time, once per deployment and
never per prediction. It refuses a \st{FAILED} model and, by default, a \st{WARNING} one, and it refuses
pickle-based formats. It reads the bytes of the model once, hashes them again and loads the model from those
same bytes, so there is no gap between the check and the use. It runs a smoke test on held-out data, computes a
\emph{release root} (the Merkle root over the manifest digests of the model and of all its ancestors) and signs
a deployment record that also binds the head of the trust log it verified against. One 32-byte value therefore
identifies the whole supply chain of a release.

A model from an outside publisher enters through a separate path. The publisher signs a manifest with its own
key, an administrator registers the publisher's \emph{public} key out of band, and the import refuses anything
unsigned or malformed, stores the rest and verifies it. A failed import is kept as quarantined evidence that
cannot be used as a parent or deployed. Its provenance ends, by design, at the identity of the publisher.

\subsection{Interoperability}
Any attestation exports as an in-toto Statement inside a DSSE envelope \cite{intotoattest,dsse}, signed with the
role key that signed the original manifest. The predicate is either the full \sys\ manifest or a SLSA
Provenance statement \cite{slsa11} in which the training or fine-tuning stage is the build, the dataset and the
base model are the resolved dependencies and the signing role is the builder. The export is verified by an
independent DSSE verifier in the test suite, including the test vector of the specification. An enterprise that
already consumes in-toto or SLSA attestations can therefore read ours.

\subsection{Design decisions and alternatives}
Table~\ref{tab:choices} justified the main choices against the state of the art. Four more decisions shaped
the implementation, each against a simpler alternative:
\begin{itemize}
\item \textbf{The manifests on disk are the source of truth and SQLite is only an index}, instead of a
  database that holds the truth. CI and auditors verify from files, and tampering with the database cannot
  fake integrity.
\item \textbf{Every use verifies the whole ancestry again}, instead of trusting the stored status of a
  parent. A stored status can be edited, and recomputing costs a few milliseconds.
\item \textbf{Trust and revocation are records of a signed log}, instead of flags in the database. Whoever
  can write a database can clear a flag (variant E5-C4).
\item \textbf{Verification happens at deployment}, instead of at each prediction. Integrity is a property of
  what is loaded, and the latency of a prediction stays unchanged.
\end{itemize}
